\usepackage{enumitem}
\usepackage{graphicx}
\usepackage{multicol}
\usepackage{xargs}
\usepackage{tikz-cd}
\usepackage{amssymb}
\usepackage{vwcol}

\usepackage{tikz}
\usetikzlibrary{arrows,calc,patterns,positioning,shapes}
\usetikzlibrary{decorations.pathmorphing}
\tikzset{
    modal/.style={>=stealth',shorten >=1pt,shorten <=1pt,auto,
    node distance=1.5cm,semithick},
    world/.style={circle,draw,minimum size=.8cm,fill=gray!15},
    point/.style={circle,draw,fill=black,inner sep=0.5mm},
    reflexive/.style={->,in=120,out=60,loop,looseness=#1},
    reflexive/.default={5},
    reflexive point/.style={->,in=135,out=45,loop,looseness=#1},
    reflexive point/.default={25},
}

\newcommand{\valentin}[1]{{\color{violet}[[V: #1]]}}

%%% Math Euler (fonts) %%%

% \DeclareMathAlphabet\EuRoman{U}{eur}{m}{n}
% \SetMathAlphabet\EuRoman{bold}{U}{eur}{b}{n}
% \newcommand{\matheu}{\EuRoman}

%%% Kripke Models

\newcommand{\props}{\mathsf{Prop}}

\newcommand{\W}{\mathrm{Q}}
\newcommand{\R}{\mathbin{\mathrm{R}}}
\newcommand{\V}{\mathrm{V}}

\newcommand{\Q}{\mathrm{Q}}

\newcommand{\Wof}[1]{\mathrm{Q}_{#1}}
\newcommand{\Qof}[1]{\mathrm{Q}_{#1}}
\newcommand{\Rof}[1]{\mathbin{\mathrm{R}_{#1}}}
\newcommand{\iRof}[1]{\mathbin{\mathrm{R}^{-1}_{#1}}}
\newcommand{\Vof}[1]{\mathrm{V}_{#1}}

\newcommand{\kmodel}{\mathfrak{K}}
\newcommand{\ikmodel}{\mathfrak{K}^{-1}}

\newcommand{\Wkm}[1][]{\Wof{\kmodel_{#1}}}
\newcommand{\Qkm}[1][]{\Qof{\kmodel_{#1}}}
\newcommand{\Rkm}[1][]{\Rof{\kmodel_{#1}}}
\newcommand{\iRkm}[1][]{\iRof{\kmodel_{#1}}}
\newcommand{\Vkm}[1][]{\Vof{\kmodel_{#1}}}

%%%% Frames

% \newcommand{\kframe}{\matheu{F}}

% \newcommand{\Wfm}[1][]{\Wof{\kframe_{#1}}}
% \newcommand{\Rfm}[1][]{\Rof{\kframe_{#1}}}
% \newcommand{\Vfm}[1][]{\Vof{\kframe_{#1}}}

% \newcommand{\fmofkm}{\kframe(\kmodel)}
% \newcommandx{\kmoffmv}[2][1={\kframe},2={\V}]{\kmodel(\npair {#1} {#2})}

%%%% Filtrations

\newcommand{\fset}{\Sigma}
\newcommandx{\fksym}[2][1=\fset,2=\kmodel]{^{#1}_{#2}}
\newcommandx{\feqr}[2][1=\kmodel,2=\fset]{\mathrel{{\equiv}{\fksym[#1][#2]}}}
\newcommandx{\feqc}[3][2=\kmodel,3=\fset]{[#1]\fksym[{#3}][{#2}]}

\newcommand{\fmodel}[1][]{\kmodel_{#1}^{\fset}}

\newcommand{\Wfm}[1][]{\Wof{\kmodel_{#1}}^{\fset}}
\newcommand{\Qfm}[1][]{\Qof{\kmodel_{#1}}^{\fset}}
\newcommand{\Rfm}[1][]{\R_{\kmodel_{#1}}^{\fset}}
\newcommand{\Vfm}[1][]{\Vof{\kmodel_{#1}}^{\fset}}


%%%% Bi-relational Models

\usepackage{xargs}
\newcommandx{\bmodel}[1][1={B}]{\mathfrak{#1}}

\newcommand{\bmsymb}{\mathfrak{B}}

\newcommand{\Wbm}[1][]{\Wof{\bmsymb_{#1}}}
\newcommand{\Qbm}[1][]{\Wof{\bmsymb_{#1}}}
\newcommand{\Rbm}[1][]{\Rof{\bmsymb_{#1}}}
\newcommand{\Vbm}[1][]{\Vof{\bmsymb_{#1}}}

\newcommand{\Fsymb}{\preccurlyeq}
\newcommand{\Fbm}{{\Fsymb_\bmsymb}}

\newcommand{\kmofbm}{\kmodel(\bmodel)}
\newcommand{\bmofhm}{\bmodel(\hmodel)}
\newcommand{\bmofwm}{\bmodel(\wmodel)}

\newcommand{\bsat}{\mathrm{b}}

%%% Hierarchical Kripke Models

\newcommand{\Kof}[1]{\mathcal{K}_{#1}}
\newcommand{\Kmodels}{\Kof{}}

\newcommand{\Fof}[1]{\Fsymb_{#1}}

\newcommand{\imodel}{\mathfrak{I}}
\newcommand{\hmodel}{\mathfrak{H}}

\newcommand{\Hleqsym}{\preccurlyeq}

\newcommand{\Hleqof}[1][]{\Hleqsym_{#1}}
\newcommand{\Hleq}{\Hleqof[\hmodel]}
\newcommand{\Fhm}[1][]{\Fof{\hmodel_{#1}}}

\newcommand{\Khm}[1][]{\Kof{\hmodel_{#1}}}

%%% Weaved Kripke Models

\newcommand{\wmodel}{\mathfrak{W}}
\newcommand{\tmodel}{\mathfrak{T}}

\newcommand{\fsim}{\rightsquigarrow}
\newcommand{\BSim}{\leftrightarroweq}

\newcommand{\Sim}{\mathbin{\mathsf{S}}}
\newcommand{\Tim}{\mathbin{\mathsf{T}}}

% \usepackage{mathrsfs}
% \newcommand{\G}{\mathcal{S}}
\newcommand{\Sof}[1]{\mathcal{S}_{#1}}
\newcommand{\Sims}{\Sof{}}

\newcommand{\Kwm}[1][]{\Kof{\wmodel_{#1}}}
\newcommand{\Swm}[1][]{\Sof{\wmodel_{#1}}}
\newcommand{\Fwm}[1][]{{\Fsymb_{\wmodel_{#1}}}}
\newcommand{\Qwm}[1][]{\Qof{\wmodel_{#1}}}

\newcommand{\Ktm}[1][]{\Kof{\tmodel_{#1}}}
\newcommand{\Stm}[1][]{\Sof{\tmodel_{#1}}}
\newcommand{\Ftm}[1][]{{\Fsymb_{\tmodel_{#1}}}}
\newcommand{\Qtm}[1][]{\Qof{\tmodel_{#1}}}

\newcommand{\fwmodel}{\wmodel^{\fset}}
\newcommand{\Kfwm}{\Kmodels^{\fset}}
\newcommand{\Sfwm}{\Sims^{\fset}}
\newcommand{\fSim}{\Sim^{\fset}}

\newcommand{\rwmodel}{\restr{\wmodel}{\kmodel}}
\newcommand{\Krwm}{\restr{\Kmodels}{\kmodel}}
\newcommand{\Srwm}{\restr{\Sims}{\kmodel}}

\newcommand{\F}{\mathcal{F}}
\newcommand{\uwmodel}{\restr{\wmodel}{\kmodel^{*}}}
\newcommand{\Kuwm}{\restr{\Kmodels}{\kmodel^{*}}}
\newcommand{\Suwm}{\restr{\Sims}{\kmodel^{*}}}

\newcommand{\pwmodel}{\restr{\wmodel}{\kmodel^{n}}}
\newcommand{\Kpwm}{\restr{\Kmodels}{\kmodel^{n}}}
\newcommand{\Spwm}{\restr{\Sims}{\kmodel^{n}}}


\newcommand{\branches}{\mathcal{C}}
\newcommand{\bwmodel}{\mathfrak{C}}
\newcommand{\Kbwm}[1][]{\Kof{\bwmodel_{#1}}}
\newcommand{\Sbwm}[1][]{\Sof{\bwmodel_{#1}}}

\newcommand{\twmodel}{\mathfrak{T}}
\newcommand{\Ktwm}{\Kof{\twmodel}}
\newcommand{\Stwm}{\Sof{\twmodel}}

\newcommand{\iwmodel}{\wmodel^{\cong}}
\newcommand{\Kiwm}{\Kmodels^{\cong}}
\newcommand{\Siwm}{\Sims^{\cong}}
\newcommand{\iSim}{\Sim^{\cong}}
\newcommand{\IsoKiwm}{\mathcal{I}^{\cong}}

\newcommand{\cwmodel}{\wmodel^{\Downarrow}}
\newcommand{\Kcwm}{\Kmodels^{\Downarrow}}
\newcommand{\Scwm}{\Sims^{\Downarrow}}

%%% Weaved Frames

% \newcommand{\wframe}{\matheu{W}}

% \newcommand{\Kwf}[1][]{\mathbb{F}_{\wframe_{#1}}}
% \newcommand{\Swf}[1][]{\Sof{\wframe_{#1}}}
% \newcommand{\Vwf}[1][]{\V_{#1}}

% \newcommand{\Qwf}[1][]{\Qof{\wframe_{#1}}}
% \newcommand{\Fwf}{{\Fsymb_\wframe}}
% \newcommand{\nFwf}{{{\not\Fsymb}_\wframe}}

%%%% From ... 2 Weaved ...

\newcommand{\QQs}[1][]{\mathcal{Q}_{#1}}
\newcommand{\Ind}{I}
\newcommand{\bigq}[1][]{\bm{Q_{#1}}}

\newcommand{\wmofhm}{\wmodel(\hmodel)}
\newcommand{\wmofbm}[1]{\wmodel_{#1}(\bmodel)}
\newcommand{\wmofwfv}{\wmodel(\wframe, \Vwf)}
\newcommand{\wfofwm}{\wmodel(\Vwf)}

%%% MISC

\usepackage{mathtools}
\DeclarePairedDelimiterX\setof[2]{\{}{\}}{\,#1 \;\delimsize\vert\; #2\,}

\DeclareMathOperator{\dom}{dom}
\DeclareMathOperator{\cod}{cod}

\newcommand{\last}[1]{#1_{{-1}}}

\newcommand{\imap}{\iota}
\newcommand{\isym}{\hookrightarrow}  % inclusion function
\newcommand{\bsym}{%
    % \hookrightarrow\mathrel{\mspace{-15mu}}\rightarrow
    % \mathbin{\tikz{\draw[>->>] (0,0) -- (0.4,0);}}
    \cong
} % bijection function

\DeclareMathOperator{\ideg}{deg}
\DeclareMathOperator{\size}{size}

\newcommand{\restr}[2]{#1\vert_{#2}}

% \DeclareMathOperator{\id}{id}

\newcommand{\bb}[1]{\mathbf{#1}}

%%% REFERENCING

\usepackage{hyperref}
\usepackage{cleveref}

\Crefname{definition}{Def.}{Defs.}
\Crefname{figure}{Fig.}{Figs.}
\Crefname{proposition}{Prop.}{Props.}
\Crefname{equation}{Eq.}{Eqs.}
\Crefname{property}{Prop.}{Props.}
\Crefname{theorem}{Thm.}{Thms.}
\Crefname{thm}{Thm.}{Thms.}
\Crefname{example}{Ex.}{Exs.}
\Crefname{section}{Sec.}{Secs.}
\Crefname{fact}{Fact}{Facts}
\Crefname{appendix}{App.}{App.}
\Crefname{table}{Tab.}{Tabs.}

%%% INLINE EQUATIONS %%%
\newcounter{ieqnnum}
\makeatletter
\newcommand*{\ieqn}[2][]{%
    \begingroup
        % Put \refstepcounter at the beginning, because
        % package `hyperref' sets the anchor here.
        \refstepcounter{ieqnnum}%
        (\roman{ieqnnum})\xspace%
        \ifx\\#1\\%
            % \label[ieqn]{\theieqnnum}
        \else
            \label[ieqn]{#1}
        \fi
        % prevent line breaks inside equation
        \relpenalty=10000 %
        \binoppenalty=10000 %
        \ensuremath{%
            % \displaystyle % larger fractions, ...
            #2%
        }%
    \endgroup
}
\makeatother

\newcounter{rieqnnum}%
\newcommand{\ieqnref}[1]{%
    \setcounter{rieqnnum}{#1}%
    \setcounter{rieqnnum}{%
        \numexpr\value{ieqnnum}+\value{rieqnnum}%
    }%
    (\roman{rieqnnum})%
}